% 三次抛硬币的自然滤过。Omega={H,T}^3 的八个等可能结果按字典序排成一行；
% 第 k 行把 Omega 按前 k 次结果分块，块数 2^k，对应 F_k 的原子；
% 块内数字是随机游走 S_k（H=+1, T=-1）在该块上的取值，S_k 在 F_k 的每块上为常数。
\begin{tikzpicture}[line width=.6pt,font=\fontsize{8}{10}\selectfont,>=Stealth,
  cell/.style={draw=FigureStroke,minimum height=6mm,inner sep=0pt},
  lab/.style={FigureInk,anchor=east},
  note/.style={FigureMuted,font=\fontsize{7.5}{9}\selectfont,anchor=west}]
% level 0
\node[cell,minimum width=80mm,fill=FigureYellowFill] at (40mm,0) {$S_0=0$};
% level 1
\foreach \i/\v in {0/{+1},1/{-1}}{
  \node[cell,minimum width=40mm,fill=FigureGreenFill] at (40mm*\i+20mm,-9mm) {$\v$};}
% level 2
\foreach \i/\v in {0/{+2},1/{0},2/{0},3/{-2}}{
  \node[cell,minimum width=20mm,fill=FigureBlueFill] at (20mm*\i+10mm,-18mm) {$\v$};}
% level 3
\foreach \i/\v in {0/{+3},1/{+1},2/{+1},3/{-1},4/{+1},5/{-1},6/{-1},7/{-3}}{
  \node[cell,minimum width=10mm,fill=FigurePurpleFill] at (10mm*\i+5mm,-27mm) {$\v$};}
% outcomes
\foreach \i/\w in {0/HHH,1/HHT,2/HTH,3/HTT,4/THH,5/THT,6/TTH,7/TTT}{
  \node[FigureMuted,font=\fontsize{7.5}{9}\selectfont] at (10mm*\i+5mm,-33.5mm) {\w};}
\node[lab] at (-2mm,0) {$\mathcal F_0$};
\node[lab] at (-2mm,-9mm) {$\mathcal F_1$};
\node[lab] at (-2mm,-18mm) {$\mathcal F_2$};
\node[lab] at (-2mm,-27mm) {$\mathcal F_3$};
\node[lab] at (-2mm,-33.5mm) {$\omega$};
\node[note] at (82mm,0) {尚无记录};
\node[note] at (82mm,-9mm) {已知第1次};
\node[note] at (82mm,-18mm) {已知前2次};
\node[note] at (82mm,-27mm) {已知全部3次};
% refinement arrows
\foreach \x in {20,60}{\draw[->,FigureGreen] (\x mm,-3.3mm) -- (\x mm,-5.7mm);}
\foreach \x in {10,30,50,70}{\draw[->,FigureBlue] (\x mm,-12.3mm) -- (\x mm,-14.7mm);}
\foreach \x in {5,15,25,35,45,55,65,75}{\draw[->,FigurePurple] (\x mm,-21.3mm) -- (\x mm,-23.7mm);}
\end{tikzpicture}
